% -----------------------------------------------------------------------
% CSC_43M04_EP — What Happens Next? Video Action Recognition
% CVPR-style double-column report
% Team 8 — GUILLAUMEY, Florian & SIGNORETTI, Andrea
% -----------------------------------------------------------------------
\documentclass[10pt,twocolumn,letterpaper]{article}

% ---- Core packages ----
\usepackage{times}
\usepackage[
  top=1in,
  bottom=1.125in,
  left=0.875in,
  right=0.875in,
  columnsep=0.375in
]{geometry}
\usepackage{graphicx}
\usepackage{amsmath,amssymb}
\usepackage{booktabs}
\usepackage{array}
\usepackage{multirow}
\usepackage{xcolor}
\usepackage{hyperref}
\usepackage{microtype}
\usepackage{enumitem}
\usepackage{caption}
\usepackage{subcaption}
\usepackage[numbers,sort&compress]{natbib}

% ---- CVPR-style title block ----
\makeatletter
\def\@maketitle{%
  \newpage\null
  \begin{center}%
    {\LARGE\bf \@title \par}%
    \vskip 1em%
    {\large \@author}%
  \end{center}%
  \vskip 1.5em}
\makeatother

% ---- Section heading sizes ----
\usepackage{titlesec}
\titleformat{\section}{\large\bfseries}{\thesection}{1em}{}
\titleformat{\subsection}{\normalsize\bfseries}{\thesubsection}{1em}{}
\titleformat{\subsubsection}{\normalsize\itshape}{\thesubsubsection}{1em}{}
\titlespacing*{\section}{0pt}{6pt}{3pt}
\titlespacing*{\subsection}{0pt}{4pt}{2pt}

% ---- Abstract environment ----
\renewenvironment{abstract}{%
  \centerline{\large\textbf{Abstract}}\vspace{0.5em}\small}{}

% ---- Handy shorthands ----
\newcommand{\pp}{\,pp}
\definecolor{leakyred}{RGB}{180,30,30}
\definecolor{cleangreen}{RGB}{30,130,60}

% -----------------------------------------------------------------------
\title{\vspace{-0.4cm}What Happens Next?\\
  Video Action Recognition on Something-Something v2}

\author{%
  \textbf{Florian GUILLAUMEY}\textsuperscript{1}\quad
  \textbf{Andrea SIGNORETTI}\textsuperscript{1}\\[2pt]
  \textsuperscript{1}École Polytechnique — CSC\_43M04\_EP\\[2pt]
  \small Team 8\quad|\quad
  \texttt{florian.guillaumey@polytechnique.edu}\quad
  \texttt{andrea.signoretti@polytechnique.edu}
}

% -----------------------------------------------------------------------
\begin{document}
\maketitle
\thispagestyle{empty}

% -----------------------------------------------------------------------
\begin{abstract}
We present our approach to the \emph{What Happens Next?} Kaggle challenge:
33-class video action recognition on a curated subset of
Something-Something v2 (SSv2)~\cite{goyal2017ssv2}, where only the first
$\approx$40\% of each clip is provided.  Track A (Closed World) achieves
44.03\% top-1 accuracy on the official validation set using a
from-scratch ensemble of Temporal Shift Module~\cite{lin2019tsm}
(TSM-ResNet18) and VideoFormer-Lite models.  Track B (Open World) reaches
\textbf{0.6586 Kaggle top-1} through a 14-model uniform ensemble built on
V-JEPA~2~\cite{bardes2025vjepa2} and VideoMAE~\cite{tong2022videomae}
backbones, with progressive unfreezing, layer-wise learning-rate decay
(LLRD), and window-capped source-frame re-extraction.  A self-discovered
data-leakage event (a discarded 0.81 submission) is documented alongside
the integrity audit that triggered a full pipeline rebuild.
\end{abstract}

% -----------------------------------------------------------------------
\section{Problem Description}
\label{sec:problem}

\paragraph{Task.}
The challenge asks models to classify 4-frame JPEG clips into one of 33
motion-defined action classes (e.g.\ \emph{Pouring something into something},
\emph{Pulling something from left to right}).  Crucially, each clip contains
only the \emph{anticipatory} portion of the original SSv2 video: the
shipped \texttt{frame\_003} sits at a median fraction of 0.40 of the
source clip, so the action's outcome is deliberately withheld.  This
design forces models to reason about motion trajectories rather than
final-state configurations.

\paragraph{Dataset.}
The provided split contains 44{,}993 training clips, 6{,}745 validation
clips, and 6{,}913 test clips across 33 classes (numeric prefixes 000–032;
class 027 is absent from the training set).  Class sizes are heavily
imbalanced: the largest class (009 — \emph{Moving something up}) has
3{,}170 training examples while the smallest (026 — \emph{Spilling
something next to something}) has only 162, a ratio of 19.6$\times$.
The KL divergence between the training and validation class distributions
is 0.088, confirming that the splits are representative.

\paragraph{Evaluation.}
Submissions are evaluated by Top-1 accuracy on a public+private Kaggle
leaderboard.  The Top-5 metric is reported internally.  A provided
baseline achieves approximately 0.55 on the private leaderboard.  The
second-place public reference score is 0.74.

\paragraph{Two-track setup.}
Track A (Closed World) prohibits pretrained models and external data;
models are trained from scratch on the provided 4-frame clips.  Track B
(Open World) permits pretrained models and external datasets provided
that privileged information (i.e.\ future frames or the ground-truth
SSv2 labels for the test set) is not used.

\paragraph{Contributions.}
Andrea Signoretti designed and implemented the full Track A pipeline
(Section~\ref{sec:trackA}), including TSM-ResNet, VideoFormer-Lite,
regularization warm-up, and the ensemble/TTA framework.  Florian
Guillaumey designed and implemented the full Track B pipeline
(Section~\ref{sec:trackB}), including source-frame re-extraction with
window-capping, the V-JEPA wrapper with progressive unfreezing and LLRD,
pseudo-labelling, the cross-preprocessing ensemble, and the data-integrity
audit.

% -----------------------------------------------------------------------
\section{Methods}
\label{sec:methods}

% ---- TRACK A ----
\subsection{Track A — Closed World}
\label{sec:trackA}

The Track A pipeline evolved through six architectural stages, guided by
ablation results on an internal training-split validation.

\subsubsection{Architecture family}

\textbf{CNN Baseline / CNN-LSTM.}  A ResNet-18 backbone with temporal
average-pooling served as the starting point.  Adding a unidirectional
LSTM over per-frame features yielded only marginal improvement due to
instability of backpropagation through time on short 4-frame sequences.

\textbf{VideoFormer-Lite (VFL).}  Per-frame ResNet-18 features are passed
to a $[\text{CLS}]$ token together with learnable temporal positional
embeddings and processed by $N$ pre-norm Transformer encoder blocks
(multi-head self-attention + MLP).  The final $[\text{CLS}]$
representation is classified.  The \emph{ultra} variant uses $N{=}4$
layers and a tuned learning rate, reaching 53.37\% on the internal
split.

\textbf{TSM-ResNet.}  The Temporal Shift Module~\cite{lin2019tsm} wraps
every residual block of ResNet-18, shifting $1/4$ of channels backward
and $1/4$ forward in time at zero parameter cost.  TSM was designed
specifically for SSv2-style motion-defined labels, making it a strong
prior for this task.

\subsubsection{Critical frame-count fix}

The single largest accuracy gain in Track A (+4.4\pp) came from aligning
the model's \texttt{num\_frames} parameter with the true clip length.
Clips contain exactly 4 real JPEG frames; setting \texttt{num\_frames\,>\,4}
triggers frame interpolation, which corrupts the temporal-shift signal
that TSM relies on.  Setting \texttt{num\_frames\,=\,4} removed this
corruption and brought \texttt{tsm\_ultra\_v2} to 57.73\% internal validation.

\subsubsection{Regularization recipe}

All Track A models share a common regularization stack whose most critical
element is a \emph{warm-up delay} on MixUp~\cite{zhang2018mixup} and
CutMix~\cite{yun2019cutmix}: these techniques are disabled for the first
10 epochs so the backbone can form basic per-frame representations before
label-mixing is applied.  This single change produced a +36\pp jump on
VFL (9.45\% $\to$ 45.78\%).  The full recipe also includes label
smoothing ($\epsilon{=}0.1$)~\cite{szegedy2016inception}, RandAugment,
random erasing~\cite{zhong2020erasing}, stochastic depth
(DropPath~0.1), dropout~0.2 before the classifier, and AdamW with cosine
annealing.  Focal loss ($\gamma{=}2$) is used in \texttt{tsm\_ultra\_v2}
to focus gradient on confusable class pairs.

CutMix probability is deliberately kept low ($p{=}0.25$) because a cut
patch spans all 4 frames and frequently occludes the hand/object motion
that defines the label.

Domain-aware augmentation is essential: horizontal flip is disabled during
training because SSv2 direction-encoded classes (e.g.\ \emph{Pulling left
to right}) change their ground-truth label under a flip.

\subsubsection{Final Track A ensemble}

The final submission fuses three checkpoints via log-softmax late fusion
with 10-crop spatial TTA (5 spatial crops $\times$ 2 flips):
\texttt{tsm\_ultra\_v2} (57.73\% internal),
\texttt{tsm\_ultra\_v2\_rotating} (rotating-fold variant trained on all
data),
and \texttt{video\_former\_lite\_ultra\_rotating}.
The ensemble reaches \textbf{44.03\% on the official \texttt{val\_dir}}.

% ---- TRACK B ----
\subsection{Track B — Open World}
\label{sec:trackB}

\subsubsection{Backbone selection}

Top-performing teams on SSv2 benchmarks consistently employ video
transformers pretrained on large unlabelled corpora.  We selected two
families:

\textbf{VideoMAE}~\cite{tong2022videomae} uses masked autoencoders on
video: 90\% of video patches are masked and the encoder must reconstruct
them, learning rich spatio-temporal representations.  We use the
HuggingFace checkpoints pretrained on Kinetics-400 (\texttt{videomae-base-finetuned-kinetics},
\texttt{videomae-large-finetuned-kinetics}).

\textbf{V-JEPA 2}~\cite{bardes2025vjepa2} is a joint-embedding predictive
architecture that predicts representations of masked video regions in
\emph{feature space} rather than pixel space.  This avoids the
reconstruction of visually uninformative regions (backgrounds, textures)
and focuses capacity on semantically relevant motion.  We use the
self-supervised checkpoint \texttt{vjepa2-vitl-fpc64-256}, which was
pretrained on unlabelled video without any SSv2 labels — a critical
requirement for integrity (Section~\ref{sec:integrity}).

\subsubsection{V-JEPA fine-tuning strategy}

V-JEPA was not present in the original codebase; we implemented the
full fine-tuning wrapper, including the frame-extraction pipeline and
the LLRD helper, from scratch.  The strategy proceeds in three stages:

\textbf{Stage 0 — Frozen attentive probe.}  A randomly initialised
linear head is trained for 6 epochs with all encoder weights frozen.
The EMA validation accuracy on the training split reaches 0.672.

\textbf{Stage 1 — Partial unfreezing with LLRD.}  The top-6 transformer
blocks are unfrozen.  Layer-wise learning-rate decay~\cite{clark2020electra,bao2022beit,he2022mae}
is applied: the head receives the base learning rate $\eta$; each lower
block $\ell$ receives $\eta \cdot \lambda^\ell$ with $\lambda{=}0.75$.
Warm-starting from Stage~0 weights, 4 epochs bring EMA validation to
0.705.

\textbf{Stage 2 — Full encoder fine-tuning with LLRD.}  All encoder
layers are unfrozen; LR decay is extended to the full depth.
Warm-starting from Stage~1, 4 more epochs bring EMA validation to 0.728.
On the official \texttt{val\_dir} evaluation, the final V-JEPA-ft
checkpoint achieves 0.6411 top-1.

The progressive unfreezing strategy follows the intuition that randomly
initialised classification heads can corrupt lower-layer representations
if allowed to propagate large gradients from the first epoch.  Freezing
the backbone initially stabilises the head, and graduated unfreezing
preserves the pretrained feature geometry.

\subsubsection{Source-frame re-extraction}
\label{sec:extraction}

V-JEPA benefits from 16 input frames, whereas only 4 frames per clip are
provided.  The folder names in the dataset encode SSv2 video IDs, so the
full source WebM files (220{,}847 videos, $\approx$19\,GB) are accessible.

We implemented \texttt{extract\_ssv2\_frames.py} to re-extract 16
real frames per clip in a \emph{window-capped} manner:
(1)~the shipped \texttt{frame\_003} is matched against low-resolution
source frames using pixel-level mean absolute error to locate the
end of the provided temporal window;
(2)~16 frames are sampled uniformly within $[0,\,\text{window\_end}]$;
(3)~the search is hard-bounded to the first 50\% of the source clip,
so the withheld action outcome ($>$60\%) is unreachable by construction.
Empirically, the resulting median window fraction is 0.39 and the maximum
is 0.50.

\subsubsection{Pseudo-labelling}
\label{sec:pseudo}

After Stage~2 fine-tuning, the V-JEPA-ft top-3 snapshot ensemble with
10-view TTA generated soft labels for the 6{,}913 test clips.  Clips
with maximum softmax probability $\geq$0.85 (100\% coverage at this
threshold) were incorporated as pseudo-labelled examples.  The model was
then continue-finetuned on
$\text{train} \cup \text{pseudo-test}$ with full encoder + LLRD for 4
more epochs, following the noisy-student
recipe~\cite{lee2013pseudo,xie2020noisy}.  The resulting V-JEPA-pseudo
achieves 0.6519 \texttt{val\_dir} and 0.6410 Kaggle; the widened
val/Kaggle gap (Section~\ref{sec:calibration}) indicated over-fitting to
the validation distribution, and pseudo-labelling was not iterated
further.

\subsubsection{Honest cross-preprocessing ensemble}
\label{sec:ensemble}

The ensemble aggregates 14 model checkpoints from three architectural
families, each fine-tuned with its own preprocessing pipeline:
\begin{itemize}[noitemsep,topsep=2pt,leftmargin=12pt]
  \item \textbf{V-JEPA} (6 snapshots): 16-frame, 256\,px, window-capped;
  \item \textbf{VideoMAE-Base K400} (3 snapshots): 4-frame, 224\,px, shipped frames;
  \item \textbf{VideoMAE-Large K400} (3 snapshots): 4-frame, 224\,px, shipped frames;
  \item \textbf{TSM-ResNet50} (2 snapshots): 4-frame, 224\,px, shipped frames.
\end{itemize}
Inference runs each model on its own preprocessing chain, producing
per-model softmax tensors that are cached to disk.
\texttt{honest\_ensemble.py} then combines them by uniform averaging of
the raw softmax scores.

The key empirical finding is that \emph{uniform averaging of diverse
honest models outperforms learned or manual weighting}.  When gradient
descent was used to fit per-model weights on validation NLL, the
optimiser concentrated $\approx$82\% of the weight on the two strongest
VideoMAE snapshots and effectively zeroed the K400 and TSM members.
Yet uniform weighting with those same weak members included scored
higher, because the architecturally distinct models produce uncorrelated
errors that average out.

\subsubsection{Infrastructure}
\label{sec:infra}

Training ran on a single RTX~4000~Ada (21\,GB VRAM).  Three
infrastructure challenges required original engineering:
(1)~a silent NFS quota (30\,GB) was producing 0-byte checkpoints until
\texttt{fsync}-durable atomic saves turned the silent failure into an
honest \texttt{OSError}, prompting migration of all checkpoints to a
local NVMe (\texttt{/Data}, 561\,GB);
(2)~a top-$K$ snapshot rename race condition caused snapshots to
overwrite each other during training — fixed by a two-phase staging
rename;
(3)~cross-architecture ensembling was silently loading Large-model
weights into a Base graph because both save \texttt{model\_name="videomae"}
— fixed by rebuilding each model from its own stored config at inference.

% -----------------------------------------------------------------------
\section{Results}
\label{sec:results}

\subsection{Overall performance}

[Table~1: Track summary results — val accuracy and Kaggle score by
track and submission type]

Table~\ref{tab:main} reports the headline results for both tracks.

\begin{table}[!t]
\centering
\caption{Summary of Track A and Track B results.
  $\dagger$~Best honest submission.}
\label{tab:main}
\setlength{\tabcolsep}{4pt}
\small
\begin{tabular}{l>{\raggedright\arraybackslash}p{3.2cm}cc}
\toprule
\textbf{Trk} & \textbf{Method} & \textbf{val} & \textbf{Kaggle}\\
\midrule
A & TSM-ResNet18 (single) & 38.25\% & — \\
A & VFL-Ultra (single) & 32.11\% & — \\
A & 3-model ens.\ + TTA & 44.03\% & — \\
\midrule
B & V-JEPA-ft standalone & 64.11\% & 0.6361 \\
B & Ens.\ v1 (9-uniform) & — & 0.6418 \\
B & V-JEPA-pseudo & 65.19\% & 0.6410 \\
B & 12-uniform & — & 0.6546 \\
B & \textbf{14-uniform}$\dagger$ & \textbf{65.19\%} & \textbf{0.6586} \\
\bottomrule
\end{tabular}
\end{table}

\subsection{Track B ablations}

\subsubsection{Honest Kaggle progression}

[Figure~1: Bar chart of Kaggle scores over submissions, red bar for
leaky/discarded result, blue bars for clean submissions]

Table~\ref{tab:progression} shows the complete Track B Kaggle history.
The 21\pp\ gap between the discarded leaky run (0.81) and the first
honest baseline (0.6361) quantifies what the data leakage was purchasing.
Every subsequent gain was earned through legitimate architectural
improvements or ensemble diversity.

\begin{table}[!t]
\centering
\caption{Track B Kaggle submission history.
  \textcolor{leakyred}{\textbf{Red}}: discarded leaky run.
  \textcolor{cleangreen}{\textbf{Green}}: honest submissions.
  $\Delta$ is relative to V-JEPA-ft (first honest baseline).}
\label{tab:progression}
\setlength{\tabcolsep}{4pt}
\small
\begin{tabular}{lcc}
\toprule
\textbf{Submission} & \textbf{Kaggle} & \textbf{$\Delta$}\\
\midrule
\textcolor{leakyred}{Leaky V-JEPA (discarded)} & \textcolor{leakyred}{0.81} & \textcolor{leakyred}{—}\\
\midrule
V-JEPA-ft standalone (clean) & 0.6361 & —\\
Ensemble v1 (9-uniform) & 0.6418 & +0.6\pp \\
V-JEPA-pseudo standalone & 0.6410 & +0.5\pp \\
V-JEPA-only (6 snapshots) & 0.6410 & regression\\
V-JEPA-heavy (3:1 manual) & 0.6499 & +1.4\pp \\
12-uniform (+ pseudo) & 0.6546 & +1.9\pp \\
\textbf{14-uniform (+ Large)} & \textbf{0.6586} & \textbf{+2.3\pp}\\
\bottomrule
\end{tabular}
\end{table}

\subsubsection{Ensemble size and weighting}

[Figure~2: Line chart of Kaggle score vs.\ number of ensemble members,
annotated to show that dropping weak models hurt]

[Table~3: Uniform vs.\ manual weighting for same model sets]

Table~\ref{tab:ensemble_ablation} isolates the effect of ensemble size
and weighting strategy on the same model collection.

\begin{table}[!t]
\centering
\caption{Ensemble strategy ablation. All rows use the same 12 checkpoints
  except where noted.}
\label{tab:ensemble_ablation}
\small
\setlength{\tabcolsep}{4pt}
\begin{tabular}{lccc}
\toprule
\textbf{Strategy} & \textbf{$n$} & \textbf{Kaggle}\\
\midrule
V-JEPA-only, uniform & 6 & 0.6410\\
9-uniform (+ K400 + TSM) & 9 & 0.6418\\
V-JEPA-heavy (3:1 weight) & 12 & 0.6499\\
12-uniform & 12 & 0.6546\\
\textbf{14-uniform} & \textbf{14} & \textbf{0.6586}\\
\bottomrule
\end{tabular}
\end{table}

The V-JEPA-only ensemble (row 1) \emph{regresses} from the 9-model
uniform ensemble despite V-JEPA's higher individual accuracy, confirming
that the weaker VideoMAE-K400 and TSM models contribute uncorrelated
errors that the ensemble corrects.  The manual V-JEPA-heavy weighting
(row 3) falls between uniform-9 and uniform-12: domain knowledge about
model quality helps less than simply adding diverse members.

\subsubsection{Progressive unfreezing}

[Figure~3: Multi-line training curve (EMA val accuracy vs.\ epoch) with
stages in different colours, showing warm-start jumps between stages]

Figure~3 shows the V-JEPA training trajectory.  The transition from the
frozen probe to Stage~1 (top-6 unfrozen + LLRD) is characterised by a
warm-start jump: Stage~1 epoch 1 (0.676) already exceeds Stage~0's final
epoch (0.672), and the curve climbs to 0.705 by epoch 4.  Stage~2 repeats
the pattern.  The warm-start jump is a direct consequence of LLRD
preserving low-layer representations while the higher learning rate on the
head adapts the classification boundary.

\subsubsection{Calibration}
\label{sec:calibration}

[Figure~4: Three-bar grouped chart (train-split EMA val, val\_dir, Kaggle)
for V-JEPA-ft and V-JEPA-pseudo]

Table~\ref{tab:calibration} quantifies the calibration gap between the
internal validation proxy and the true Kaggle score.

\begin{table}[!t]
\centering
\caption{Calibration: train-split EMA val, official \texttt{val\_dir},
  and Kaggle, for two V-JEPA checkpoints. The pseudo-label run's
  \texttt{val\_dir}$\to$Kaggle gap doubles relative to the clean fine-tune.}
\label{tab:calibration}
\small
\setlength{\tabcolsep}{3pt}
\begin{tabular}{lccc}
\toprule
\textbf{Model} & \textbf{Train-split val} & \textbf{\texttt{val\_dir}} & \textbf{Kaggle}\\
\midrule
V-JEPA-ft & 0.7281 & 0.6411 & 0.6361\\
V-JEPA-pseudo & 0.7415 & 0.6519 & 0.6410\\
\bottomrule
\end{tabular}
\end{table}

The \texttt{val\_dir}$\to$Kaggle gap is 0.005 for V-JEPA-ft and 0.011 for
V-JEPA-pseudo — a doubling that warned us against further pseudo-label
iterations.  The train-split $\to$ Kaggle gap ($\approx$0.09) reflects
that the internal train/val split was drawn from the same distribution,
while the Kaggle test set is independently sampled.

\subsection{Track A analysis}

\subsubsection{Ensemble component contributions}

[Table~4: Leave-one-out ablation for Track A 3-model ensemble]

Table~\ref{tab:trackA_ablation} shows the Track A leave-one-out ablation on
the full \texttt{val\_dir} (6{,}746 clips).  All three components make a
positive contribution; removing \texttt{tsm\_ultra\_v2} alone costs more
than the other two combined, confirming it as the dominant member.

\begin{table}[!t]
\centering
\caption{Track A leave-one-out ablation on \texttt{val\_dir}
  (no TTA, 3-model weighted ensemble).}
\label{tab:trackA_ablation}
\small
\setlength{\tabcolsep}{3pt}
\begin{tabular}{lcc}
\toprule
\textbf{Configuration} & \textbf{Top-1} & \textbf{$\Delta$}\\
\midrule
Full ensemble (3 models) & 40.43\% & —\\
$-$ \texttt{tsm\_ultra\_v2} & 37.26\% & $-$3.17\pp\\
$-$ \texttt{tsm\_ultra} & 39.56\% & $-$0.87\pp\\
$-$ \texttt{video\_former\_lite\_ultra} & 39.82\% & $-$0.61\pp\\
\bottomrule
\end{tabular}
\end{table}

\subsubsection{Per-family accuracy (Track B)}

[Figure~5: Horizontal bar chart of single-model val\_dir top-1 accuracy by
backbone family, colour-coded, with the 11-model uniform ensemble shown
as the top bar]

Table~\ref{tab:family_acc} shows single-family \texttt{val\_dir} top-1
accuracy versus the 11-model uniform ensemble.  No individual family
reaches the ensemble score, and the TSM family --- despite its low
standalone accuracy (0.34) --- contributes to the ensemble's improvement
because its errors are the least correlated with those of any
transformer-based family (inter-family agreement 0.37–0.40;
Table~\ref{tab:diversity}).

\begin{table}[!t]
\centering
\caption{Per-family \texttt{val\_dir} top-1 accuracy (Track B, 3 snapshots
  per family, single-view) and the 11-model uniform ensemble.}
\label{tab:family_acc}
\small
\setlength{\tabcolsep}{4pt}
\begin{tabular}{lcc}
\toprule
\textbf{Family} & \textbf{Snapshots} & \textbf{val top-1}\\
\midrule
VideoMAE-Base K400 4f & 3 & 0.5576\\
VideoMAE-Large K400 4f & 2 & 0.5945\\
VideoMAE-Base SSv2 4f & 3 & 0.6047\\
TSM ResNet-50 4f & 3 & 0.3364\\
\midrule
\textbf{Uniform ensemble (11 models)} & 11 & \textbf{0.6222}\\
\bottomrule
\end{tabular}
\end{table}

\begin{table}[!t]
\centering
\caption{Pairwise inter-family prediction agreement on the test set.
  Lower agreement implies more complementary errors.}
\label{tab:diversity}
\small
\setlength{\tabcolsep}{3pt}
\begin{tabular}{llc}
\toprule
\textbf{Family A} & \textbf{Family B} & \textbf{Agreement}\\
\midrule
V-JEPA-pseudo & TSM & 0.370\\
V-JEPA-ft & TSM & 0.372\\
K400-Base & TSM & 0.375\\
V-JEPA-ft & K400-Base & 0.582\\
V-JEPA-ft & V-JEPA-pseudo & 0.854\\
\bottomrule
\end{tabular}
\end{table}

Table~\ref{tab:diversity} quantifies prediction diversity.  TSM is the
most complementary member: its pairwise agreement with every transformer
family is below 0.40, compared to 0.58 for two VideoMAE variants.  The
intra-family agreement for V-JEPA (0.85) explains why a V-JEPA-only
ensemble underperforms a diverse one of equal size.

\subsection{Failure analysis}

[Figure~6: 33$\times$33 confusion matrix heatmap for the Track B
11-model ensemble on val\_dir]

\subsubsection{Dominant error patterns}

Three confusion patterns dominate both tracks and reflect the design
intent of the SSv2 benchmark:

\textbf{Real vs.\ pretended actions.}  The worst single confusion (Track B)
is \emph{Picking something up} (class 011) being predicted as
\emph{Pretending to pick something up} (class 014): 56 errors out of
199 validation clips, an 11.6\% top-1 F1 for class 011.  The two classes
share an identical early-clip trajectory; the disambiguating cue
(whether the object actually rises) falls in the withheld future.  This
is precisely the anticipation reasoning that the task is designed to
probe.  Focal loss ($\gamma{=}2$) in Track A partially addresses this but
does not resolve it.

\textbf{Temporal direction confusion.}  \emph{Folding} $\leftrightarrow$
\emph{Unfolding} (24–29 bidirectional errors) share identical spatial
appearance; with 4 frames the temporal ordering is fragile.

\textbf{Static sink class.}  \emph{Holding something} attracts many
misclassified clips from motion classes (\emph{Moving something down}:
30 errors; \emph{Opening}: 28), acting as the network's default prediction
when no decisive motion is detected in the provided frames.

\subsubsection{Per-class analysis}

The 5 best-predicted Track A classes (F1 $>$ 0.54) are all single-arc
motions: \emph{Pulling from left to right}, \emph{Pouring}, \emph{Moving
closer}, \emph{Uncovering}, \emph{Folding}.  The 5 worst (F1 $<$ 0.20)
are \emph{Picking up}, \emph{Pretending to put}, \emph{Pretending to throw},
\emph{Spilling}, and \emph{Showing to camera} — classes where early-clip
evidence is ambiguous or intentionally misleading.

The correlation between training-set class size and per-class accuracy is
weak ($r{=}0.24$), confirming that the difficulty stems from semantic
ambiguity rather than sample scarcity.

\subsubsection{The top-5 recovery gap}

The top-5 accuracy on \texttt{val\_dir} is 0.882, compared to 0.622 top-1,
a gap of 0.260.  This gap is largest for the hardest classes:
\emph{Pretending to put something into something} recovers from 11.8\%
top-1 to 79.4\% top-5 (+67.6\pp), and \emph{Picking something up} from
11.6\% to 72.9\% (+61.3\pp).  The model assigns high probability to the
correct class but ranks a visually similar ``pretend'' variant first.
This suggests the model understands the action family but cannot reliably
distinguish real from pretended execution without observing the outcome.

\subsection{Data-integrity story}
\label{sec:integrity}

[Figure~7: Histogram of frame\_003 position (fraction of source video)
for sampled clips, with dashed line at 0.40 and shaded region $>$0.50]

An early extraction run sampled frames uniformly over the full source
clip ($0$--$100\%$).  Training with V-JEPA on these frames yielded 0.81
Kaggle --- 7\pp\ above the public second-place.  An audit revealed two
concurrent leakage vectors:
(1)~frames from the withheld outcome portion ($>$60\%) were reachable,
giving the model access to information the competition explicitly
withholds;
(2)~the initial V-JEPA checkpoint used was finetuned on the SSv2
dataset with ground-truth labels, meaning the backbone had seen the
test-set videos during its own pretraining.

Both runs were discarded and the entire pipeline was rebuilt as described
in Section~\ref{sec:extraction}, switching to the self-supervised
\texttt{vjepa2-vitl-fpc64-256} backbone and enforcing the 50\% hard
bound on extraction.  We consider the 21\pp\ honest/leaky accuracy
difference a direct measurement of the combined leakage effect.

% -----------------------------------------------------------------------
\section{Conclusion}
\label{sec:conclusion}

We have presented a two-track solution to video action anticipation on
SSv2.  Track A demonstrates that carefully engineered from-scratch models
--- with domain-appropriate regularization and correct temporal sampling
--- can reach competitive accuracy (44.03\%) on a notoriously hard
benchmark.  Track B shows that progressive unfreezing with LLRD enables
effective fine-tuning of large video transformers on a constrained GPU
budget, and that \emph{ensemble diversity across architectural families
consistently outperforms single-family ensembling or learned weighting
schemes}.

Failure analysis reveals a fundamental limit: the dataset's design intent
is to distinguish real from pretended actions and directional from
reverse motions — disambiguations that require observing the withheld
outcome.  Improving on these classes likely requires either explicit
trajectory prediction heads or iterative pseudo-label refinement with
stronger confidence calibration.

\paragraph{Future work.}
The most promising next steps are: (1) including V-JEPA-Giant once
inference-time memory constraints are addressed; (2) targeted class-weight
rebalancing for the \emph{holding} and \emph{pretending} sinks; (3) a
second round of pseudo-labelling with temperature-scaled confidence
thresholds to avoid the val-distribution over-fitting observed after the
first round; and (4) 17-model ensemble with a newly fine-tuned
VideoMAE-Large 16-frame window-capped checkpoint currently in training.

% -----------------------------------------------------------------------
{\small
\bibliographystyle{plainnat}
\bibliography{references}
}

% -----------------------------------------------------------------------
% APPENDIX: Figure and Table placeholders / data
% -----------------------------------------------------------------------
\clearpage
\appendix

\section*{Figures and Tables}

\paragraph{Figure~1 — Honest Kaggle progression.}
Vertical bar chart of Kaggle scores over all Track B submissions.
The leaky bar (0.81) is shown in red with a strike-through label
\emph{``DISCARDED''}. All honest bars are shown in blue/green.
A horizontal dashed line marks the 0.74 second-place reference.
Data source: Table~\ref{tab:progression}.

\paragraph{Figure~2 — Ensemble size vs.\ Kaggle score.}
Line + scatter plot. $x$-axis: number of models in the uniform ensemble;
$y$-axis: Kaggle top-1. An annotation arrow at $n{=}6$ marks the
V-JEPA-only regression: dropping the K400 and TSM members despite their
lower individual accuracy \emph{hurts} the ensemble.
Data source: Table~\ref{tab:ensemble_ablation}.

\paragraph{Figure~3 — V-JEPA progressive-unfreezing training curves.}
Multi-line plot of EMA validation accuracy vs.\ epoch on the training split.
Stage 0 (frozen probe), Stage 1 (top-6 unfrozen + LLRD), and Stage 2
(full encoder + LLRD) are plotted in distinct colours.
The warm-start jump at each stage transition is annotated.
Data:

\medskip
\begin{center}
\footnotesize
\setlength{\tabcolsep}{3pt}
\begin{tabular}{llc}
\toprule
Stage & Epoch & EMA val acc\\
\midrule
Frozen probe & 1 & 0.5755\\
Frozen probe & 6 & 0.6725\\
Stage 1 (top-6 + LLRD) & 1 & 0.6764\\
Stage 1 (top-6 + LLRD) & 4 & 0.7052\\
Stage 2 (full + LLRD) & 1 & 0.7015\\
Stage 2 (full + LLRD) & 4 & 0.7281\\
Pseudo retrain & 1 & 0.7180\\
Pseudo retrain & 4 & 0.7415\\
\bottomrule
\end{tabular}
\end{center}

\paragraph{Figure~4 — Calibration comparison.}
Grouped bar chart with three bars per model:
train-split EMA val, \texttt{val\_dir}, Kaggle.
Data source: Table~\ref{tab:calibration}.

\paragraph{Figure~5 — Per-family \texttt{val\_dir} accuracy.}
Horizontal bar chart sorted by accuracy, colour-coded by family
(V-JEPA, VideoMAE-Base SSv2, VideoMAE-Large, VideoMAE-Base K400, TSM).
Data source: Table~\ref{tab:family_acc}.

\paragraph{Figure~6 — Confusion matrix.}
$33\times33$ heatmap of the 11-model Track B ensemble on \texttt{val\_dir}.
Colour scale: white = 0, dark blue = max. The strongest off-diagonal
clusters correspond to: \emph{Picking up}$\to$\emph{Pretending to pick up}
(col 014 in row 011), \emph{Folding}$\leftrightarrow$\emph{Unfolding}
(rows/cols 003,\,032), \emph{Moving down/up}$\to$\emph{Holding} (col 005).
Full $33\times33$ CSV data available in supplementary material.

\paragraph{Figure~7 — Frame position histogram.}
Histogram of \texttt{frame\_003}'s fractional position in the source clip,
sampled from 14 training clips.  A vertical dashed line at 0.40 marks the
median end of the provided window.  The shaded region 0.50--1.00 is
labelled \emph{``withheld future — data leakage if extracted''}.
An overlay shows the corrected extraction's window-end distribution
(median 0.39, max 0.50).

\end{document}
